% ============================================================
%  LEMBAR SOAL  --  MATEMATIKA  --  SM ITS 2024 + Latihan
%  Paket 7.1  |  45 Soal
% ============================================================

\documentclass[11pt]{article}
\usepackage[utf8]{inputenc}
\usepackage[T1]{fontenc}
\usepackage[a4paper,margin=2.5cm,top=2.8cm]{geometry}
\usepackage{amsmath,amssymb}
\usepackage{graphicx}
\usepackage{enumitem}
\usepackage{multicol}
\usepackage{xcolor}
\usepackage[most]{tcolorbox}
\usepackage{fancyhdr}
\usepackage{booktabs}
\usepackage{icomma}
\usepackage{array}
\usepackage{tikz}

\definecolor{mapelcolor}{RGB}{20,110,60}
\definecolor{mapelbg}{RGB}{228,246,234}

\pagestyle{fancy}\fancyhf{}
\lhead{\small SM ITS -- Paket 7.1}
\rhead{\small Matematika}
\cfoot{\small\thepage}
\renewcommand{\headrulewidth}{0.4pt}

\newtcolorbox{soal}[1]{
  breakable,
  colback=white, colframe=black!70,
  coltitle=white, colbacktitle=black,
  fonttitle=\bfseries\sffamily,
  title=Nomor #1,
  boxrule=0.8pt, arc=3pt,
  left=3mm, right=3mm, top=2mm, bottom=2mm}

\newenvironment{pil}{%
  \begin{enumerate}[label=(\alph*),leftmargin=*,itemsep=1pt,topsep=3pt]}%
  {\end{enumerate}}
\newenvironment{pilB}{%
  \par\smallskip\begin{multicols}{2}
  \begin{enumerate}[label=(\alph*),leftmargin=*,itemsep=1pt,topsep=3pt]}%
  {\end{enumerate}\end{multicols}}

\linespread{1.05}

\begin{document}
\thispagestyle{empty}

\begin{center}
  \fbox{\parbox{0.6\textwidth}{\centering\bfseries\large LEMBAR SOAL}}
\end{center}
\vspace{0.5cm}
\begin{center}{\LARGE\bfseries MATEMATIKA}\end{center}
\vspace{0.5cm}

\begin{center}
  \begin{tabular}{@{\hspace{2cm}}ll@{\hspace{0.4cm}:@{\hspace{0.4cm}}}l}
    Jenjang      & & SMA / Sederajat               \\[0.3em]
    Hari/Tanggal & & \underline{\hspace{5cm}}      \\[0.3em]
    Waktu        & & \underline{\hspace{5cm}}      \\[0.3em]
    Paket Soal   & & 7.1                           \\
  \end{tabular}
\end{center}

\vspace{0.5cm}\noindent\rule{\linewidth}{0.6pt}\vspace{0.3cm}

\noindent\textbf{\underline{PETUNJUK UMUM}}
\begin{enumerate}[leftmargin=*,itemsep=4pt,topsep=4pt]
  \item Anda \textbf{tidak} diperbolehkan menggunakan sumber-sumber eksternal,
        termasuk internet, \textit{large language model} seperti ChatGPT, Claude,
        LLaMA, Gemini, Meta AI, dan sebagainya. Anda sangat dianjurkan untuk
        mencoba mengerjakan sendiri terlebih dahulu karena evaluasi ini dibuat
        untuk \textbf{mengukur} tingkat kepahaman; nilai $<70$ mengindikasikan
        \textbf{tidak lulus}, dan jika sudah melewati \textit{deadline} maka
        dianggap tidak mengerjakan yang berarti \textbf{tidak lulus}.

  \item Terdapat 2 tahap pengerjaan, yaitu tahap \textbf{Try Out} dan tahap
        \textbf{Review}. Try Out bersifat \textit{open book}; Anda diperbolehkan
        membuka buku apapun selama bersifat \textit{offline} dengan rentang waktu
        2 jam pada tanggal \underline{\hspace{3cm}}. Setelah itu memasuki
        tahap Review dengan \textit{schedule} rentang tanggal
        \underline{\hspace{3cm}}--\underline{\hspace{3cm}}, di mana Anda
        diharapkan mengerjakan ulang dengan disertakan penjelasan/pembelaan
        mengapa revisi jawaban adalah pilihan tersebut, dan tetap mencantumkan
        penjelasan pada soal yang walaupun pada penilaian sudah benar.

  \item Anda dianjurkan untuk mencantumkan caranya walaupun tidak termasuk
        penilaian, mengingat sifatnya adalah sebagai bahan pembelajaran.
        Mohon bertanggung jawab atas pekerjaan Anda sendiri.

  \item Bacalah setiap soal dengan teliti. Terdapat tiga jenis soal:
        \begin{enumerate}[label=(\arabic*),leftmargin=2em,topsep=2pt,itemsep=2pt]
          \item \textbf{Pilihan Ganda Biasa} --- 5 opsi (A--E), pilih
                1 jawaban paling tepat;
          \item \textbf{Pilihan Ganda Majemuk Kompleks} --- tabel
                \textbf{Benar}/\textbf{Salah}, tentukan kebenaran tiap
                pernyataan secara mandiri; dan
          \item \textbf{Isian Singkat} --- tulis nilai/ekspresi pada
                kolom yang tersedia.
        \end{enumerate}

  \item Soal berjumlah \textbf{45 butir} soal dengan bobot asli per soal adalah
        $\mathbf{\tfrac{20}{9}}$ untuk mendapatkan nilai sempurna sebesar \textbf{100}.

  \item Silakan bertanya apabila terdapat soal yang ambigu, rancu, atau kurang
        spesifik selama itu tidak menjurus kepada jawaban soal tersebut.

  \item Isilah detail informasi berikut.\\[0.3em]
        \textbf{Nama}\hspace{0.45cm}:\enspace\underline{\hspace{8cm}}
\end{enumerate}

\vspace{0.6cm}
\noindent\textit{Dengan menandatangani kolom di bawah ini, saya menyatakan telah membaca,
memahami, dan menyetujui seluruh peraturan yang tercantum di atas.}

\vspace{0.6cm}
\noindent\fbox{\parbox{5cm}{\vspace{2.2cm}\centering\small Tanda Tangan Peserta\vspace{0.3cm}}}

\clearpage

% ===================================================================
%  BAGIAN A: SOAL ASLI SM ITS 2024  (Q1--22)
% ===================================================================

{\noindent\large\bfseries Bagian A \textendash\ Rekonstruksi Soal Asli SM ITS 2024\par}
\smallskip\hrule\medskip

\begin{soal}{1 \normalfont\small[Matriks]}
Diketahui $\begin{pmatrix}2&1\\1&3\end{pmatrix}\begin{pmatrix}x&2\\1&y\end{pmatrix}=\begin{pmatrix}7&10\\5&20\end{pmatrix}$.
Nilai $x-y$ adalah \ldots
\begin{pilB}
  \item $-4$ \item $-3$ \item $3$ \item $4$ \item $6$
\end{pilB}
\end{soal}

\begin{soal}{2 \normalfont\small[Logaritma]}
Jika $\log_2 3=p$ dan $\log_2 5=q$, maka $\log_6 45=\ldots$
\begin{pil}
  \item $\dfrac{2p+q}{1+p}$ \item $\dfrac{p+2q}{1+p}$
  \item $\dfrac{1+2p}{p+q}$ \item $\dfrac{2+p+q}{1+q}$ \item $\dfrac{p+q}{p+1}$
\end{pil}
\end{soal}

\begin{soal}{3 \normalfont\small[Integral Substitusi]}
Hasil dari $\displaystyle\int\dfrac{6x-4}{\sqrt{3x^2-4x+7}}\,dx$ adalah \ldots
\begin{pil}
  \item $2\sqrt{3x^2-4x+7}+C$
  \item $\sqrt{3x^2-4x+7}+C$
  \item $\tfrac{1}{2}\sqrt{3x^2-4x+7}+C$
  \item $\tfrac{2}{3}\sqrt{3x^2-4x+7}+C$
  \item $3\sqrt{3x^2-4x+7}+C$
\end{pil}
\end{soal}

\begin{soal}{4 \normalfont\small[Geometri Ruang]}
Balok $10\times6\times4$ cm. Semut bergerak dari satu titik sudut ke titik sudut berhadapan melalui permukaan. Jarak terpendek adalah \ldots
\begin{pilB}
  \item $4\sqrt{17}$ cm \item $10\sqrt{2}$ cm \item $2\sqrt{58}$ cm \item $8\sqrt{5}$ cm \item $6\sqrt{7}$ cm
\end{pilB}
\end{soal}

\begin{soal}{5 \normalfont\small[Turunan Fungsi]}
$f(x)=ax^2+bx-1$ dan $g(x)=bx^2+ax+2$. Jika $f'(2)=6$ dan $g'(-2)=-7$, maka $3a-2b=\ldots$
\begin{pilB}
  \item $-3$ \item $-1$ \item $1$ \item $3$ \item $5$
\end{pilB}
\end{soal}

\begin{soal}{6 \normalfont\small[Limit Fungsi]}
$\displaystyle\lim_{x\to c}(f(x)-g(x))=4$ dan $\displaystyle\lim_{x\to c}f(x)g(x)=5$.
Nilai $\displaystyle\lim_{x\to c}(f(x)^2+g(x)^2)=\ldots$
\begin{pilB}
  \item $16$ \item $21$ \item $26$ \item $36$ \item $30$
\end{pilB}
\end{soal}

\begin{soal}{7 \normalfont\small[Vektor]}
$\vec{a}=(k,2,0)$ dan $\vec{b}=(1,0,0)$; $k>0$; sudut antara $\vec{a}$ dan $\vec{b}$ adalah $60°$.
Nilai $k=\ldots$
\begin{pil}
  \item $\dfrac{\sqrt{3}}{3}$ \item $\dfrac{2\sqrt{3}}{3}$ \item $\sqrt{3}$ \item $2\sqrt{3}$ \item $3\sqrt{3}$
\end{pil}
\end{soal}

\begin{soal}{8 \normalfont\small[Pertidaksamaan Linear]}
Nilai $x$ yang memenuhi $2(3x-1)-5\leq x+8$ adalah \ldots
\begin{pilB}
  \item $x\leq3$ \item $x\geq3$ \item $x<-3$ \item $x>-3$ \item $x\leq-3$
\end{pilB}
\end{soal}

\begin{soal}{9 \normalfont\small[Komposisi Fungsi]}
$f(x)=x^2-4x+1$ dan $g(x)=2x-3$. Maka $f(g(x))=\ldots$
\begin{pil}
  \item $4x^2-20x+22$ \item $4x^2-12x+10$
  \item $2x^2-20x+22$ \item $4x^2+20x+22$ \item $4x^2-20x+10$
\end{pil}
\end{soal}

\begin{soal}{10 \normalfont\small[Definisi Turunan]}
$f(x)=3x^3-5x^2+7x-1$. Nilai $\displaystyle\lim_{h\to0}\dfrac{f(x+h)-f(x)}{h}=\ldots$
\begin{pil}
  \item $9x^2-10x+7$ \item $3x^2-5x+7$
  \item $9x^2+10x+7$ \item $x^2-10x+7$ \item $9x^2-10x-7$
\end{pil}
\end{soal}

\begin{soal}{11 \normalfont\small[Fungsi Implisit]}
$f(3x-2)=2x^2+x+5$. Nilai $f(4)=\ldots$
\begin{pilB}
  \item $13$ \item $15$ \item $17$ \item $19$ \item $21$
\end{pilB}
\end{soal}

\begin{soal}{12 \normalfont\small[Operasi Pecahan]}
Hitunglah $\dfrac34\div\dfrac98\times\dfrac25=\ldots$
\begin{pilB}
  \item $\dfrac{4}{15}$ \item $\dfrac{5}{12}$ \item $\dfrac{2}{15}$ \item $\dfrac{8}{45}$ \item $\dfrac{1}{6}$
\end{pilB}
\end{soal}

\begin{soal}{13 \normalfont\small[Limit]}
Nilai $\displaystyle\lim_{x\to9}\dfrac{x-9}{\sqrt{x}-3}=\ldots$
\begin{pilB}
  \item $3$ \item $6$ \item $9$ \item $12$ \item $15$
\end{pilB}
\end{soal}

\begin{soal}{14 \normalfont\small[Aljabar Dasar]}
Hasil pengembangan $(x-4)(x+7)=\ldots$
\begin{pil}
  \item $x^2+3x-28$ \item $x^2-3x-28$
  \item $x^2+11x-28$ \item $x^2-11x+28$ \item $x^2-3x+28$
\end{pil}
\end{soal}

\begin{soal}{15 \normalfont\small[Luas Daerah]}
Luas daerah yang dibatasi $y=x^3$ dan $y=x$ adalah \ldots
\begin{pilB}
  \item $\dfrac14$ \item $\dfrac13$ \item $\dfrac12$ \item $1$ \item $\dfrac34$
\end{pilB}
\end{soal}

\begin{soal}{16 \normalfont\small[Nilai Mutlak]}
Himpunan penyelesaian $|2x-3|<5$ adalah \ldots
\begin{pil}
  \item $-1<x<4$ \item $x<-1$ atau $x>4$
  \item $-4<x<1$ \item $x\leq-1$ atau $x\geq4$ \item $-2<x<4$
\end{pil}
\end{soal}

\begin{soal}{17 \normalfont\small[Pecahan]}
Hitunglah $\dfrac12+\dfrac23-\dfrac34+\dfrac16=\ldots$
\begin{pilB}
  \item $\dfrac{5}{12}$ \item $\dfrac{7}{12}$ \item $\dfrac34$ \item $\dfrac{11}{12}$ \item $\dfrac{1}{12}$
\end{pilB}
\end{soal}

\begin{soal}{18 \normalfont\small[Aritmetika]}
Tiga per delapan dari $1440$ adalah \ldots
\begin{pilB}
  \item $480$ \item $520$ \item $540$ \item $560$ \item $600$
\end{pilB}
\end{soal}

\begin{soal}{19 \normalfont\small[Barisan Logaritma]}
Barisan $\ln3,\;\ln9,\;\ln27,\;\ln81,\;\ldots$ Jumlah delapan suku pertama adalah \ldots
\begin{pilB}
  \item $28\ln3$ \item $32\ln3$ \item $36\ln3$ \item $40\ln3$ \item $44\ln3$
\end{pilB}
\end{soal}

\begin{soal}{20 \normalfont\small[Deret Aritmetika]}
Jumlah dua belas bilangan kelipatan 4 pertama adalah \ldots
\begin{pilB}
  \item $288$ \item $300$ \item $312$ \item $336$ \item $360$
\end{pilB}
\end{soal}

\begin{soal}{21 \normalfont\small[Eksponen]}
Bentuk sederhana $\dfrac{8x^3y^{-2}z^5}{24x^{-1}y^4z^{-3}}=\ldots$
\begin{pil}
  \item $\tfrac13x^4y^{-6}z^8$ \item $\tfrac13x^2y^2z^2$
  \item $3x^4y^{-6}z^8$ \item $\tfrac13x^{-4}y^6z^{-8}$ \item $3x^{-4}y^6z^{-8}$
\end{pil}
\end{soal}

\begin{soal}{22 \normalfont\small[Pecahan Campuran]}
Hitunglah $\dfrac56\div\dfrac{10}{9}+\dfrac34=\ldots$
\begin{pilB}
  \item $\dfrac43$ \item $\dfrac32$ \item $\dfrac53$ \item $\dfrac74$ \item $\dfrac{11}{12}$
\end{pilB}
\end{soal}

% ===================================================================
%  BAGIAN B: SET LATIHAN 1  (Q23--37)
% ===================================================================

\medskip{\noindent\large\bfseries Bagian B \textendash\ Set Latihan 1\par}
\smallskip\hrule\medskip

\begin{soal}{23 \normalfont\small[Trigonometri]}
Diketahui $\tan a=\tfrac12$, $\tan b=\tfrac13$, dan $\tan c=\tfrac15$.
Nilai $\tan(a+b+c)=\ldots$
\begin{pilB}
  \item $\tfrac12$ \item $1$ \item $\tfrac32$ \item $2$ \item $3$
\end{pilB}
\end{soal}

\begin{soal}{24 \normalfont\small[Turunan Trigonometri]}
$y=\cos\!\left(\dfrac3x\right)$. Maka $\dfrac{dy}{dx}=\ldots$
\begin{pil}
  \item $\sin\tfrac3x$ \item $-\sin x$
  \item $\sin3x$ \item $\dfrac{3}{x^2}\sin\dfrac3x$ \item $\sin\dfrac2x$
\end{pil}
\end{soal}

\begin{soal}{25 \normalfont\small[Persamaan Kuadrat]}
$x^2+(m-1)x+9$ mempunyai akar-akar real berbeda. Batasan nilai $m$ adalah \ldots
\begin{pil}
  \item $-5<m<7$ \item $m<-5$ atau $m>7$
  \item $m<-7$ atau $m>5$ \item $-7<m<5$ \item $m>-7$ atau $m<5$
\end{pil}
\end{soal}

\begin{soal}{26 \normalfont\small[SPLDV]}
Toko buku: 2 buku gambar $+$ 8 buku tulis $=$ Rp48.000; 3 gambar $+$ 5 tulis $=$ Rp37.000.
Harga 1 buku gambar dan 2 buku tulis adalah \ldots
\begin{pil}
  \item Rp24.000 \item Rp20.000 \item Rp17.000 \item Rp13.000 \item Rp14.000
\end{pil}
\end{soal}

\begin{soal}{27 \normalfont\small[Turunan Fungsi]}
$f(x)=x^3+3x^2-9x-7$ turun pada interval \ldots
\begin{pil}
  \item $1<x<3$ \item $-1<x<3$
  \item $-3<x<1$ \item $3<x<1$ \item $5<x<3$
\end{pil}
\end{soal}

\begin{soal}{28 \normalfont\small[Kombinatorika]}
Banyak bilangan terdiri dari angka berlainan antara 100 dan 400 dari angka $1,2,3,4,5$ adalah \ldots
\begin{pilB}
  \item $36$ \item $48$ \item $52$ \item $60$ \item $68$
\end{pilB}
\end{soal}

\begin{soal}{29 \normalfont\small[Lingkaran]}
Lingkaran berpusat $P(3,-1)$ melalui $(5,2)$. Persamaannya adalah \ldots
\begin{pil}
  \item $x^2+y^2+6x-2y-55=0$ \item $x^2+y^2+6x-2y-31=0$
  \item $x^2+y^2-6x+2y-23=0$ \item $x^2+y^2-6x+2y-3=0$
  \item $x^2+y^2+6x+2y-13=0$
\end{pil}
\end{soal}

\begin{soal}{30 \normalfont\small[Kombinatorika]}
Banyak cara perjalanan $M\to O\to M$ via $N$ (4 jalur $M\!-\!N$, 5 jalur $N\!-\!O$), tidak boleh melalui jalur yang sama saat pulang, adalah \ldots
\begin{pilB}
  \item $180$ \item $240$ \item $260$ \item $160$ \item $120$
\end{pilB}
\end{soal}

\begin{soal}{31 \normalfont\small[Titik Berat 3D]}
$A(1,0,2)$, $B(5,4,10)$, $C(0,-1,6)$. Titik berat segitiga $ABC$ adalah \ldots
\begin{pil}
  \item $(2,-1,6)$ \item $(2,1,6)$ \item $(3,-1,6)$ \item $(3,2,6)$ \item $(6,1,3)$
\end{pil}
\end{soal}

\begin{soal}{32 \normalfont\small[Trigonometri]}
$BD=CD$, $AB=p$ satuan, $\angle BDA=x$. Panjang $BC=\ldots$
\begin{pil}
  \item $\dfrac{p}{\tan x}$ \item $2p\tan x$ \item $p\tan x$ \item $\dfrac{2p}{\tan x}$ \item $2p\sin x$
\end{pil}
\end{soal}

\begin{soal}{33 \normalfont\small[Vektor]}
$\vec a$ dan $\vec b$ membentuk sudut tumpul $\alpha$ dengan $\sin\alpha=\tfrac{1}{\sqrt7}$; $|\vec a|=\sqrt5$, $|\vec b|=\sqrt7$. Nilai $\vec a\cdot\vec b=\ldots$
\begin{pilB}
  \item $30$ \item $\sqrt{30}$ \item $-\sqrt{30}$ \item $-20$ \item $-30$
\end{pilB}
\end{soal}

\begin{soal}{34 \normalfont\small[Limit Trigonometri]}
$\displaystyle\lim_{x\to0}\dfrac{\tan x}{x^2\csc3x}=\ldots$
\begin{pilB}
  \item $1$ \item $2$ \item $3$ \item $4$ \item $5$
\end{pilB}
\end{soal}

\begin{soal}{35 \normalfont\small[Parabola dan Garis]}
Parabola $y=x^2$ memotong garis $y=x+2$ di $A$ dan $B$. Panjang $AB=\ldots$
\begin{pilB}
  \item $4$ \item $3$ \item $2$ \item $2\sqrt2$ \item $3\sqrt2$
\end{pilB}
\end{soal}

\begin{soal}{36 \normalfont\small[Permutasi]}
${}^6P_3+{}^5P_2+{}^7P_4=\ldots$
\begin{pilB}
  \item $900$ \item $980$ \item $600$ \item $680$ \item $390$
\end{pilB}
\end{soal}

\begin{soal}{37 \normalfont\small[Deret Geometri Tak Hingga]}
Deret geometri tak hingga, suku awal $a$ dan rasio $r$; $a+r=6$ dan $S_\infty=5$. Nilai $\dfrac{a}{r}=\ldots$
\begin{pilB}
  \item $-20$ \item $25$ \item $\dfrac56$ \item $-\dfrac{1}{25}$ \item $-25$
\end{pilB}
\end{soal}

% ===================================================================
%  BAGIAN C: SOAL TANTANGAN  (Q38--45)
%  Setiap soal merupakan kolaborasi 3 bab
% ===================================================================

\medskip{\noindent\large\bfseries Bagian C \textendash\ Soal Tantangan (Kolaborasi 3 Bab)\par}
\smallskip\hrule\medskip

\begin{soal}{38 \normalfont\small[Turunan $\times$ Integral $\times$ Geometri Analitik]}
Garis singgung kurva $y=x^3-3x^2+2$ di titik $(2,-2)$ membentuk daerah tertutup dengan kurva.
Luas daerah tersebut adalah \ldots
\begin{pilB}
  \item $\dfrac{9}{4}$ \item $\dfrac{27}{4}$ \item $\dfrac{9}{2}$ \item $\dfrac{27}{2}$ \item $9$
\end{pilB}
\end{soal}

\begin{soal}{39 \normalfont\small[Trigonometri $\times$ Persamaan $\times$ Barisan]}
Semua nilai $x\in[0,\pi]$ yang memenuhi $\cos3x+\sin3x=1$ membentuk suatu himpunan.
Jumlah semua nilai $x$ tersebut adalah \ldots
\begin{pilB}
  \item $\pi$ \item $\dfrac{4\pi}{3}$ \item $\dfrac{3\pi}{2}$ \item $\dfrac{5\pi}{3}$ \item $2\pi$
\end{pilB}
\end{soal}

\begin{soal}{40 \normalfont\small[Vektor $\times$ Geometri Analitik $\times$ Sistem Persamaan]}
Diketahui $A(1,2)$, $B(5,4)$, dan $C(p,q)$ membentuk segitiga siku-siku di $C$
(yaitu $\vec{CA}\perp\vec{CB}$) dengan luas $\triangle ABC=5$.
Jumlah semua nilai $p+q$ yang mungkin adalah \ldots
\begin{pilB}
  \item $8$ \item $10$ \item $12$ \item $14$ \item $16$
\end{pilB}
\end{soal}

\begin{soal}{41 \normalfont\small[Barisan Geometri $\times$ Eksponen $\times$ Logaritma]}
Tiga bilangan positif membentuk barisan geometri dengan jumlah 7 dan hasil kali 8.
Jika barisan bersifat menaik, nilai ${}^2\!\log$ dari suku ke-4 adalah \ldots
\begin{pilB}
  \item $2$ \item $3$ \item $4$ \item $6$ \item $8$
\end{pilB}
\end{soal}

\begin{soal}{42 \normalfont\small[Matriks $\times$ Sistem Persamaan $\times$ Aljabar]}
Diketahui $A=\begin{pmatrix}a&b\\b&a\end{pmatrix}$ dan $A^2=\begin{pmatrix}13&12\\12&13\end{pmatrix}$
dengan $a>b$. Nilai $a^2-b^2+ab=\ldots$
\begin{pilB}
  \item $7$ \item $9$ \item $11$ \item $13$ \item $15$
\end{pilB}
\end{soal}

\begin{soal}{43 \normalfont\small[Eksponen $\times$ Identitas Aljabar $\times$ Persamaan]}
Jika $4^x+4^{-x}=14$, maka $2^x+2^{-x}=\ldots$
\begin{pilB}
  \item $2$ \item $3$ \item $4$ \item $2\sqrt{3}$ \item $2\sqrt{5}$
\end{pilB}
\end{soal}

\begin{soal}{44 \normalfont\small[Logaritma $\times$ Sistem Persamaan $\times$ Fungsi]}
Pasangan $(x,y)$ memenuhi sistem $\begin{cases}\log_2(xy)=5\\\log_2\dfrac{x}{y}=1\end{cases}$ ($x,y>0$).
Nilai $x+y=\ldots$
\begin{pilB}
  \item $8$ \item $10$ \item $12$ \item $14$ \item $16$
\end{pilB}
\end{soal}

\begin{soal}{45 \normalfont\small[Limit $\times$ Faktorisasi Polinomial $\times$ Aljabar]}
Nilai $\displaystyle\lim_{x\to2}\frac{x^4-3x^3+2x^2+x-2}{x^3-4x^2+5x-2}=\ldots$
\begin{pilB}
  \item $3$ \item $4$ \item $5$ \item $6$ \item $7$
\end{pilB}
\end{soal}

\end{document}
